\chapter{Artefakt 2 - Automatisierung von Logistics Lines}
\label{chap:Art2LL}

Dieses Kapitel stellt die \emph{Choreografiebasierte Automatisierung und MTP-ba\-sier\-te Integration von Logistics Lines} als zweites Artefakt dieser Dissertation vor. Dabei werden bestehende Choreografiekonzepte aus \cite{Stutz.2026,Stutz.2021} übernommen und um neu spezifizierte MTP-Konzepte ergänzt. Die Nutzung von Choreografiekonzepten für die Automatisierung von Logistics Lines wurde im Rahmen der studentischen Arbeit \cite{Ortmann.2021} untersucht und in Teilen in \cite{Blumenstein.ETFA,Blumenstein.MoProLog,Stutz.2022} veröffentlicht. Die Evaluationsbeispiele des \emph{BASF-Demonstrators} und der \emph{Steuerungstechnischen Emulation} in Abschnitt~\ref{sec:Evaluationsbeispiele} veranschaulichen die Anwendung des Artefakts auf industrienahe Anwendungsfälle der produktionsnahen Logistik. Darüber hinaus enthält das GitHub-Repository zu dieser Dissertation \cite{Blumenstein.Github} Detailspezifikationen der eingeführten MTP-Erweiterungen sowie eine noch detailliertere Beschreibung ihrer Anwendung auf die Automatisierung einer Logistics Line zur Sackabfüllung.

\section{Artefaktübersicht}
\label{sec:Art2LLUebersicht}

Für die Automatisierung von Logistics Lines wird gemäß Abschnitt~\ref{subsec:AnforderungenLL} ein choreografiebasierter Ansatz verfolgt, der auf den in \cite{Stutz.2021,Stutz.2026} beschriebenen Mechanismen basiert. Diese werden auf die produktionsnahe Logistik übertragen und um MTP-basierte Schnittstellen sowie Modelle zur automatisierten Integration erweitert. Dabei wird zwischen der horizontalen Integration von LEAs innerhalb einer Logistics Line und der vertikalen Integration choreografierter Logistics Lines in einen übergeordneten LOL unterschieden.

\pagebreak

\textbf{Horizontale Integration von LEAs in eine Logistics Line:} Die horizontalen Interaktionen zwischen LEAs werden mittels eines Choreografie-Konfigurators im LOL konfiguriert, sodass eine choreografierte Logistics Line entsteht. Zur automatisierten Integration der LEAs in den Choreografie-Konfigurator werden in den LEA-MTPs deren verschaltbare Ein- und Ausgänge in Form \textbf{semantischer MTP-Modelle} beschrieben (Abschnitt~\ref{subsec:LeaChoreoModelle}). Auf dieser Basis wird eine Choreografiekonfiguration entworfen und anschließend auf die LEA-Steuerungen übertragen. Für diese Übertragung sind \textbf{MTP-basierte Konfigurationsschnittstellen} in den LEA-Steuerungen vorgesehen und in den LEA-MTPs beschrieben (Abschnitt~\ref{subsec:LeaChoreoSchnittstellen}). Abschnitt~\ref{subsec:ChoreoSet} führt die notwendigen MTP-basierten Schnittstellen und Modelldefinitionen in einem neuen MTP-Aspekt \emph{ChoreographySet} ein.

\textbf{Vertikale Integration von Logistics Lines:} Eine choreografierte Logistics Line als Ganzes kann in einen übergeordneten LOL integriert und automatisiert durch ein Auftrags-Management oder manuell durch einen Operator gesteuert werden. Hierfür stellt sie eine gemeinsame Diensteschnittstelle (\textbf{Line-Schnittstelle}, Abschnitt~\ref{subsec:LineSchnittstelle}) bereit, die die choreografierte Logistics Line integrierbar und ansteuerbar macht, als wäre sie ein einzelner MTP-Dienst. Zudem wird ein Bedienbild zur Bedienung der Logistics Line durch einen Operator bereitgestellt (\textbf{Line-Bedienbild}, Abschnitt~\ref{subsec:LinienBedienbild}). Die Line-Schnittstelle sowie das Line-Bedienbild werden in einem gemeinsamen MTP für die Logistics Line (\textbf{Composed-MTP}, Abschnitt~\ref{subsec:KonzeptComposedMTP}) beschrieben. Abschnitt~\ref{sec:ChoreoVertikal} beschreibt eine Methodik, wie die Schnittstellen einzelner LEAs zu einer Line-Schnittstelle und die Bedienbilder einzelner LEAs zu einem Line-Bedienbild zusammengesetzt werden können. Zudem wird die Modellierung von Line-Schnittstellen und Line-Bedienbildern in Form eines Composed-MTP eingeführt.

\section{Horizontale Integration von Logistics Equipment Assemblies in eine Logistics Line}
\label{sec:ChoreoHorizontal}

Dieser Abschnitt beschreibt die horizontale Integration von LEAs in eine choreografierte Logistics Line. Einführend stellt Abschnitt~\ref{subsec:EinfuehrungsBeispielLL} ein veranschaulichendes Beispiel einer choreografierten Logistics Line vor und verdeutlicht die horizontalen Interaktionen zwischen den LEAs. Anschließend folgen die Konzepte zur technischen Umsetzung. In den Abschnitten~\ref{subsec:LineArchitektur} und \ref{subsec:ChoreoLeaArchitektur} wird die Architektur choreografierter Logistics Lines und der enthaltenen LEAs in Anlehnung an \cite{Stutz.2026} dargestellt. Abschnitte~\ref{subsec:LeaChoreoModelle} und \ref{subsec:LeaChoreoSchnittstellen} führen daraufhin notwendige MTP-basierte Modelle und Schnittstellen zur Choreografiekonfiguration von Logistics Lines ein, die im Rahmen dieser Dissertation entwickelt wurden. 

\subsection{Anwendungsbeispiel}
\label{subsec:EinfuehrungsBeispielLL}

Abbildung~\ref{fig:EinfuehrungsBeispielLL} zeigt zur Veranschaulichung eine beispielhafte Logistics Line zur Sackabfüllung und -palettierung bestehend aus drei LEAs -- einer Abfüllmaschine (engl. \gls{ffs}), einem Transportband (engl. \gls{conv}) und einem Lagenpalettierer (engl. Palletizer, PAL). Zur horizontalen Steuerung der Logistics Line gemäß dem Choreografieprinzip sind parametrierende, prozedurale, verriegelnde und regelnde Relationen zwischen den LEAs sowie Konstanten definiert. Die Relationen beschreiben, welche Informationen die LEAs von den anderen LEAs beobachten und wie sie auf diese Informationen reagieren.

\begin{figure}[htbp]
	\centerline{\includegraphics[width=1.0\textwidth]{Inhalt/Abbildungen/07_Linen-Automatisierung/Beispiel_Linie_Interaktionen.pdf}}
	\caption{Anwendungsbeispiel einer choreografierten Logistics Line mit drei Logistics Equipment Assemblies zur Veranschaulichung der horizontalen Interaktionen}
	\label{fig:EinfuehrungsBeispielLL}
\end{figure}

Beispielsweise ist ein geordnetes Anfahren der Line von hinten nach vorn, in der Reihenfolge PAL -- CONV -- FFS, umgesetzt. Hierbei werden notwendige Parameter, wie z.~B. die zu verpackende Sackanzahl, am PAL eingestellt und durch \textbf{parametrierende Relationen} der FFS zur Verfügung gestellt (Abbildung~\ref{fig:EinfuehrungsBeispielLL}, (1)). Das bedeutet, die FFS beobachtet die Parameter des PAL und überträgt deren Werte auf ihre eigenen Parameter. Mittels \textbf{prozeduraler Relationen} erfolgt die Steuerung der Dienstezustände der LEAs. Wurde beispielsweise der Dienst des PAL gestartet, startet als nächstes der Dienst des CONV (Abbildung~\ref{fig:EinfuehrungsBeispielLL}, (2)). Hierzu beobachtet der CONV den Dienstezustand des PAL. Sobald dieser den Zustand EXECUTE erreicht, startet der CONV seinen eigenen Dienst. Neben direkten Relationen, bei denen die Änderung genau einer beobachteten Variable eine Aktion in der beobachtenden LEA auslöst, sind auch komplexere Logiken in den Relationen denkbar. So gehen im Fehlerfall (z.~B. STOPPED-Zustand eines LEA-Dienstes) auch alle anderen LEA-Dienste in den STOPPED-Zustand. Für z.~B. die PAL-LEA bedeutet dies, dass sie die Zustände des FFS- und des CONV-Dienstes beobachtet. Sobald einer der beiden Dienste in den STOPPED-Zustand wechselt (Logisches ODER), versetzt die PAL-LEA auch ihren eigenen Dienst in STOPPED. Während des Betriebs der Line werden zwischen den LEAs kontinuierlich Freigabesignale übertragen, die als MTP-Prozesswerte von einer nachgelagerten auf eine vorgelagerte LEA übertragen werden. Sie wirken als \textbf{verriegelnde Relationen} und stellen sicher, dass ein LO erst dann in eine LEA einfahren darf, wenn diese dazu bereit ist (Abbildung~\ref{fig:EinfuehrungsBeispielLL}, (3)). So beobachtet beispielsweise der CONV den Freigabesignalausgang des PAL und überträgt dessen Wert auf seinen eigenen Freigabesignaleingang. Die Logik des CONV sorgt wiederum dafür, dass eine LO-Übergabe nur bei gesetztem Freigabesignal erfolgt. Zudem sind \textbf{regelnde Relationen} denkbar, mittels derer beispielsweise die Arbeitsgeschwindigkeit der LEAs einer Logistics Line aneinander angepasst werden kann. Schließlich können auch \textbf{Konstanten} choreografiebasiert gesetzt werden. Diese werden beispielsweise genutzt, um die Betriebsarten der LEA-Dienste zu definieren, oder um konstante Parameter zu setzen. Regelnde Relationen und Konstanten sind im oben gezeigten einfachen Beispiel nicht dargestellt. Eine detaillierte Beschreibung dieses Anwendungsbeispiels ist im GitHub-Repository zu dieser Dissertation \cite{Blumenstein.Github} zu finden. 

\subsection{Systemarchitektur einer choreografierten Logistics Line}
\label{subsec:LineArchitektur}

	% - Allgemeine Abläufe innerhalb einer Choreografie
	% - Systemarchitektur einer Linie --> Einführung Aktiver Choreografie-Teilnehmer

Die technische Umsetzung des in Abschnitt~\ref{subsec:EinfuehrungsBeispielLL} beschriebenen Prinzips erfolgt auf Basis des Choreografie-Konzepts gemäß \cite{Stutz.2021,Stutz.2026}. Abbildung~\ref{fig:ChoreoArchitektur} stellt die grundlegende Architektur einer choreografierten Logistics Line dar.

\begin{figure}[htbp]
	\centerline{\includegraphics[width=0.6\textwidth]{Inhalt/Abbildungen/07_Linen-Automatisierung/Choreo_Abstrakte_Architektur.pdf}}
	\caption{Abstrakte Systemarchitektur einer choreografierten Logistics Line in Anlehnung an \cite{Stutz.2021}}
	\label{fig:ChoreoArchitektur}
\end{figure}

Das Steuerungssystem jeder LEA einer Logistics Line stellt den Ausführungskontext für einen MTP-Dienst sowie die notwendige Basisfunktionalität hinsichtlich Datenverarbeitung und Kommunikation zur Teilnahme an einer Choreografie bereit (in \cite{Stutz.2026} \textbf{Active Choreography Participant} genannt). In jeder LEA-Steuerung ist ein Dienst gemäß Kapitel~\ref{chap:Art1LEA} vom LEA-Hersteller implementiert, welcher das \textbf{Native Steuerungsprogramm} gemäß \cite{Stutz.2026} darstellt. Um die in Abschnitt~\ref{subsec:EinfuehrungsBeispielLL} beschriebenen Relationen zwischen den LEAs umzusetzen, wird dieser Dienst befähigt, Variablen der Dienste der anderen LEAs innerhalb der Choreografie zu beobachten und sich entsprechend deren Werten zu verhalten. In \cite{Stutz.2026} wird zwischen internen (innerhalb einer LEA) und externen Relationen (zwischen verschiedenen LEAs) unterschieden. Zur Umsetzung der Interaktionen zwischen den LEAs einer Logistics Line sind ausschließlich externe Relationen notwendig, die über konfigurierbare Kommunikationsmechanismen umgesetzt sind (in \cite{Stutz.2026} \textbf{Konfigurierbare Kommunikation} genannt). Jede LEA beinhaltet zudem intern konfigurierte Verhaltensregeln, die beschreiben, wie die LEA auf die ihr zur Verfügung gestellten Signale anderer LEAs reagieren soll (in \cite{Stutz.2026} \textbf{Konfigurierbare Logik} genannt). 

\subsection{Architektur einer LEA als Choreografieteilnehmer}
\label{subsec:ChoreoLeaArchitektur}

	% - Architektur eines Choreografie-Teilnehmers --> Einführung Configurable Logic und Configurable Communication
	% - Komponenten der Configurable Logic
	% - Komponenten der Configurable Communication und UnionType

Abbildung~\ref{fig:ParticipantArchitektur} zeigt die grundlegende Architektur eines Active Choreography Participant (im Folgenden \textbf{Choreografieteilnehmer} genannt) gemäß \cite{Stutz.2026}. Dieser Architektur muss eine LEA folgen, um an einer Choreografie teilnehmen zu können.

\begin{figure}[htbp]
	\centerline{\includegraphics[width=0.85\textwidth]{Inhalt/Abbildungen/07_Linen-Automatisierung/Architektur_ChoreoParticipant_2.pdf}}
	\caption{Architektur einer LEA als Teilnehmer einer Choreografie in Anlehnung an \cite{Stutz.2026} }
	\label{fig:ParticipantArchitektur}
\end{figure}
 
Jede LEA-Steuerung enthält einen LEA-Dienst gemäß Kapitel~\ref{chap:Art1LEA}, der durch weitere Softwarekomponenten befähigt wird, an einer Choreografie teilzunehmen. Hierfür sieht \cite{Stutz.2026} die Entwurfsmuster \textbf{Configurable Logic} und \textbf{Configurable Communication} zur Umsetzung der Konzepte der Konfigurierbaren Logik und Konfigurierbaren Kommunikation vor. 

\subsubsection{Configurable Logic}

Die Komponenten der Configurable Logic (Abbildung~\ref{fig:ParticipantArchitektur}, grüne Färbung) dienen der konfigurierbaren Verarbeitung der zur Verfügung stehenden Informationen, welche LEA-intern oder LEA-extern bereitgestellt werden können. Ausgaben des LEA-Dienstes sind dabei LEA-interne Eingaben der Konfigurierbaren Logik (verwaltet in einer \textbf{Input-Liste}). Zudem werden Informationen anderer LEAs als LEA-externe Eingaben in der Input-Liste verwaltet. Die Informationen der Input-Liste können durch konfigurierbare Funktionen verarbeitet werden, die in einer \textbf{Logic-Liste} verwaltet werden. Die resultierenden Ausgaben der Konfigurierbaren Logik (verwaltet in einer \textbf{Output-Liste}) werden LEA-intern als Eingaben an den LEA-Dienst zurückgegeben und können dort Zustandsübergänge bzw. Änderungen an Prozesswerten, Parametern oder Reportwerten hervorrufen. Zudem können in der Output-Liste verwaltete Informationen anderen LEAs als LEA-externe Eingaben zur Verfügung gestellt werden. Abbildung~\ref{fig:KonfigLogikKonzept} stellt die Bestandteile des Entwurfsmusters Configurable Logic gemäß \cite{Stutz.2026} dar.

\begin{figure}[htbp]
	\centerline{\includegraphics[width=1.0\textwidth]{Inhalt/Abbildungen/07_Linen-Automatisierung/Software_Pattern_Konfigurierbare_Logik.drawio.pdf}}
	\caption{Bestandteile des Entwurfsmusters Configurable Logic gemäß \cite{Stutz.2026}}
	\label{fig:KonfigLogikKonzept}
\end{figure}

Im Mittelpunkt steht eine Komponente \emph{ConfigurableLogic}, welche die Steuerung der Ausführung der Konfigurierbaren Logik übernimmt. Hierzu hat sie Zugriff auf die Input-, Logic- und Output-Liste. Während der Ausführung der \emph{ConfigurableLogic} werden Werte aus der Input-Liste gelesen, Verarbeitungen vorgenommen und die Resultate in die Output-Liste geschrieben. In \cite{Stutz.2026} werden zudem notwendige Datentypen und Enumerationen vorgestellt, die das in den drei Listen abgelegte Informationsmodell der Konfigurierbaren Logik beschreiben. Eine genauere Vorstellung dieses Informationsmodells und der Arbeitsweise der Komponente \emph{ConfigurableLogic} ist im Anhang~\ref{subsec:AnhangConfigurableLogic} zu finden. 

\subsubsection{Configurable Communication}
 
Die Komponenten der Configurable Communication (Abbildung~\ref{fig:ParticipantArchitektur}, blaue Färbung) dienen dem konfigurierbaren Informationsaustausch zwischen den LEAs innerhalb einer Choreografie. Ziel ist, dass jede LEA die LEA-externen Informationen erhält, die sie für die Ausführung der in ihrer Logic-Liste konfigurierten Verarbeitungsfunktionen benötigt. Um die Kommunikation zwischen LEAs verschiedener Hersteller zu ermöglichen, soll eine standardisierte Kommunikationstechnologie genutzt werden. Im Einklang mit den bisherigen Konzepten im MTP-Umfeld wird im Rahmen dieser Dissertation die Nutzung von OPC~UA Client/Server vorgesehen. Für diese Kommunikationstechnologie sieht \cite{Stutz.2026} sowohl die Kommunikationsvariante des aktiven Lesens als auch des aktiven Schreibens von Informationen vor. Im Falle des aktiven Lesens werden Informationen aus der Input-Liste einer anderen LEA ausgelesen und in die eigene Input-Liste übertragen. Im Falle des aktiven Schreibens werden Informationen aus der eigenen Output-Liste in die Input-Liste einer anderen LEA geschrieben. Zudem kann das aktive Schreiben gemäß \cite{Stutz.2026,Stutz.2023} verwendet werden, um nicht choreografiebefähigtes Equipment mittels dezentraler Orchestrierung anzubinden. Um alle diese Varianten zu ermöglichen, werden im Rahmen dieser Dissertation MTP-Konzepte sowohl für das aktive Lesen als auch Schreiben vorgesehen. Wie in \cite{Stutz.2026} analysiert, wird allerdings die Variante des aktiven Lesens empfohlen. Somit wird diese Variante besonders fokussiert und auch im Rahmen der Evaluation umgesetzt sowie bewertet. Abbildung~\ref{fig:KonfigKommunikationKonzept} stellt die Komponenten des Entwurfsmusters Configurable Communication gemäß \cite{Stutz.2026} dar.

\begin{figure}[htbp]
	\centerline{\includegraphics[width=1.0\textwidth]{Inhalt/Abbildungen/07_Linen-Automatisierung/Software_Pattern_Konfigurierbare_Kommunikation_Übersicht.drawio.pdf}}
	\caption{Bestandteile des Entwurfsmusters Configurable Communication ge\-mäß \cite{Stutz.2026}}
	\label{fig:KonfigKommunikationKonzept}
\end{figure}

Im Mittelpunkt steht die Komponente \emph{OpcUaClientServerManager}, die die Verwaltung und Ausführung der Konfigurierbaren Kommunikation übernimmt. Der \emph{OpcUaClientServerManager} verwaltet Komponenten zur Konfiguration und Aus\-füh\-rung von OPC~UA-Verbindungen (\emph{UaConnections}) zwischen den Teilnehmern sowie, je nach Kommunikationsvariante, Komponenten zur Konfiguration und Ausführung von Lese- und/oder Schreiboperationen (\emph{UaReader} bzw. \emph{UaWriter}). Zudem hat der \emph{OpcUaClientServerManager} Zugriff auf die Input- und Output-Listen. Für die Kommunikationsvariante des aktiven Lesens werden durch konfigurierte \emph{UaReader} Werte aus anderen LEAs gelesen und in die eigene Input-Liste der LEA übertragen. Im Fall des aktiven Schreibens existiert zudem eine \emph{ValueList} als passiver Eingangskanal. Auf diesen kann durch entsprechende \emph{UaWriter} von anderen LEAs geschrieben werden. Diese Werte werden anschließend auf bestimmte Indizes der Input-Liste übertragen. Stutz \cite{Stutz.2026} stellt notwendige Datentypen und Enumerationen vor, die das Informationsmodell der Konfigurierbaren Kommunikation beschreiben. Eine genauere Vorstellung dieses Informationsmodells und der Arbeitsweise der Komponenten ist im Anhang~\ref{subsec:AnhangConfigurableCommunication} zu finden. 

\subsubsection{Union Type}
\label{subsec:UnionType}

Die zwischen den LEAs einer Logistics Line ausgetauschten Informationen weisen gemäß den Konzepten in \cite{Stutz.2026} den in den Abbildungen~\ref{fig:UnionType} gezeigten Datentyp \emph{UnionType} auf. 

\begin{figure}[htbp]
	\centerline{\includegraphics[width=0.35\textwidth]{Inhalt/Abbildungen/07_Linen-Automatisierung/UnionType.drawio.pdf}}
	\caption{Datentyp UnionType gemäß \cite{Stutz.2026} }
	\label{fig:UnionType}
\end{figure}

\pagebreak

Die Besonderheit dieses Typs ist, dass er (aktuell fünf) Variablen verschiedener Datentypen vorhält, aus denen eine zur Nutzung ausgewählt wird. Dies ermöglicht die Festlegung des Datentyps zur Laufzeit auch bei Steuerungssystemen, die eine Definition des Datentyps bereits zum Engineeringzeitpunkt erfordern.


\subsection{Semantische Modelle zur LEA-Integration in einen Choreografie-Konfigurator}
\label{subsec:LeaChoreoModelle}

	% - Horizontale Integration bedarf ein Konfigurationssystem im LOL
	% - Bedarf der im folgenden vorgestellten Bestandteile
	% - Integration von LEAs in einen Choreografie-Konfigurator --> Input- und Output-Liste
	% - Interne Inputs/outputs
	% - Externe Inputs/outputs

Für einen Choreografie-Konfigurator im LOL ist es notwendig, den Inhalt der Input- und Output-Liste einer LEA semantisch eindeutig identifizieren zu können. Abbildung~\ref{fig:InOutputListen} stellt den grundlegenden Aufbau dieser Listen schematisch dar. 

\begin{figure}[htbp]
	\centerline{\includegraphics[width=0.6\textwidth]{Inhalt/Abbildungen/07_Linen-Automatisierung/InOutputListen.pdf}}
	\caption{Aufbau der Input- und Output-Listen einer LEA}
	\label{fig:InOutputListen}
\end{figure}

Innerhalb dieser Listen lassen sich grundsätzlich fixe, konfigurierbare und schreibbare Elemente unterscheiden. Die \textbf{fixen Elemente} sind durch die LEA-Pro\-gram\-mie\-rung fest vorgegeben und nicht veränderbar. Im Falle der Input-Liste handelt es sich um fixe Ausgaben des LEA-Dienstes (z.~B. dessen Zustand oder aktuelle Prozedur), im Falle der Output-Liste um fixe Eingaben an den LEA-Dienst (z.~B. Befehle an den Dienst oder Prozedurvorgabe). Bei den \textbf{konfigurierbaren Elementen} handelt es sich um Platzhalter für Werte, die aus anderen LEAs gelesen (im Falle der Input-Liste) oder an andere LEAs geschrieben (im Falle der Output-Liste) werden sollen. \textbf{Schreibbare Elemente} sind Platzhalter für Werte, die aus anderen LEAs in die Input-Liste geschrieben werden. Beim LEA-Engineering wird eine fixe Anzahl an konfigurierbaren und schreibbaren Elementen festgelegt, die im Zuge der Choreografiekonfiguration mit konkreten Relationen zu anderen LEAs verbunden werden können.

Als Basis für die Abbildung dieser Inhalte in einem MTP wird die Modelldefinition \emph{ChoreographyParticipant} eingeführt, die eine LEA als möglichen Choreografieteilnehmer ausweist. Zur Abbildung der Input- und Output-Listen der LEA werden die MTP-Modelldefinitionen \emph{InputList} und \emph{OutputList} eingeführt, welche entsprechende \emph{InputElements} bzw. \emph{OutputElements} organisieren. Gemäß der Unterscheidung nach fixen, konfigurierbaren und schreibbaren Elementen werden die in Tabelle~\ref{tab:InOutputElements} gezeigten Modelldefinitionen konkreter Input- und Output-Elemente eingeführt.\footnote{Die zugehörigen Schnittstellendefinitionen werden im nachfolgenden Abschnitt~\ref{subsec:LeaChoreoSchnittstellen} beschrieben.}

% Modelldefinitionen InputElements und OutputElements
\begin{footnotesize}
	\begin{longtable}{|
		>{\columncolor[HTML]{E0E0E0}}l
		>{\columncolor[HTML]{E0E0E0}}l 
		>{\columncolor[HTML]{FFFFFF}}l |}
		\caption{Arten von Elementen in der Input- und Output-Liste einer LEA}
		\label{tab:InOutputElements}\\
		\hline

		\multicolumn{1}{|L{2.5cm}|}{\cellcolor[HTML]{E0E0E0}\textbf{Modell-definition}}        	& 
		\multicolumn{1}{L{2.6cm}|}{\cellcolor[HTML]{E0E0E0}\textbf{Schnittstellen-definition}} &
		\multicolumn{1}{L{9.35cm}|}{\cellcolor[HTML]{E0E0E0}\textbf{Beschreibung}} 			
		\\ \hline

		\multicolumn{1}{|L{2.5cm}|}{\cellcolor[HTML]{FFFFFF} FixedInput-Element}                     & 
			\multicolumn{1}{L{2.6cm}|}{\cellcolor[HTML]{FFFFFF} UnionElement}   &
		\multicolumn{1}{L{9.35cm}|}{\cellcolor[HTML]{FFFFFF} Diese Definitionen charakterisieren fest definierte Ausgaben eines LEA-Dienstes, die als Eingaben für die Konfigurierbare Logik dienen (z.~B. Service States). Diese Elemente der Input-Liste sind durch Konfiguration nicht veränderbar. }    
		\\ \hline

		\multicolumn{1}{|L{2.5cm}|}{\cellcolor[HTML]{FFFFFF} FixedOutput-Element}                     & 
		\multicolumn{1}{L{2.6cm}|}{\cellcolor[HTML]{FFFFFF} UnionElement}    &
		\multicolumn{1}{L{9.35cm}|}{\cellcolor[HTML]{FFFFFF} Diese Definitionen charakterisieren fest definierte Ausgaben der Konfigurierbaren Logik, die als Eingaben für den LEA-Dienst dienen (z.~B. Service Commands). Diese Elemente der Output-Liste sind durch Konfiguration nicht veränderbar.}          
		\\ \hline

		\multicolumn{1}{|L{2.5cm}|}{\cellcolor[HTML]{FFFFFF} Configurable-InputElement}                     & 
		\multicolumn{1}{L{2.6cm}|}{\cellcolor[HTML]{FFFFFF} UnionElement}     &  
		\multicolumn{1}{L{9.35cm}|}{\cellcolor[HTML]{FFFFFF} Diese Definitionen charakterisieren konfigurierbare Eingaben für die Konfigurierbare Logik, welche extern von anderen LEAs der Logistics Line bereitgestellt werden und über Leseoperationen eingelesen werden (relevant für die Kommunikationsvariante des aktiven Lesens). Durch Konfiguration wird festgelegt, von wo externe Informationen eingelesen werden.}                
		\\ \hline

		\multicolumn{1}{|L{2.5cm}|}{\cellcolor[HTML]{FFFFFF} Configurable-OutputElement}                     & 
		\multicolumn{1}{L{2.6cm}|}{\cellcolor[HTML]{FFFFFF} UnionElement} &
		\multicolumn{1}{L{9.35cm}|}{\cellcolor[HTML]{FFFFFF} Diese Definitionen charakterisieren konfigurierbare Ausgaben der Konfigurierbaren Logik, welche durch Schreiboperationen extern anderen LEAs der Logistics Line bereitgestellt werden (relevant für die Kommunikationsvariante des aktiven Schreibens). Durch Konfiguration wird festgelegt, wohin die Informationen extern geschrieben werden.}              
		\\ \hline

		\multicolumn{1}{|L{2.5cm}|}{\cellcolor[HTML]{FFFFFF} WritableInput-Element}                     & 
		\multicolumn{1}{L{2.6cm}|}{\cellcolor[HTML]{FFFFFF} WritableUnion-Element}  &
		\multicolumn{1}{L{9.35cm}|}{\cellcolor[HTML]{FFFFFF} Diese Definitionen charakterisieren passive Eingaben für die Konfigurierbare Logik, auf die von anderen LEAs der Logistics Line Informationen geschrieben werden können (relevant für die Kommunikationsvariante des aktiven Schreibens). }              
		\\ \hline

	\end{longtable}
\end{footnotesize}

\subsection{Schnittstellen zur Konfiguration der horizontalen Interaktion}
\label{subsec:LeaChoreoSchnittstellen}

	% - Konfigurierbaren Logik --> ChoreographyParticipantManager-Schnittstelle
	% - Konfigurierbaren Kommunikation --> OpcUaClientServerManager-Schnittstelle

Zum Herunterladen einer Choreografiekonfiguration aus einem Choreografie-Kon\-fi\-gu\-ra\-tor in die LEA-Steuerung sind MTP-basierte Schnittstellen zur Konfiguration der Konfigurierbaren Logik und der Konfigurierbaren Kommunikation notwendig.

\subsubsection{Konfigurierbare Logik}

Abschnitt~\ref{subsec:ChoreoLeaArchitektur} stellt die Komponente \emph{ConfigurableLogic} zur Steuerung der Aus\-füh\-rung der Konfigurierbaren Logik gemäß \cite{Stutz.2026} vor. Zur Interaktion mit dieser Komponente wird in dieser Dissertation die MTP-Schnitt\-stel\-len\-de\-fi\-ni\-ti\-on \emph{ChoreographyParticipantManager}\footnote{Die \emph{ChoreographyParticipantManager}-Schnittstelle kann auch als dynamisches Bedienbildelement in das LEA-Bedienbild eingebunden werden.} spezifiziert. Sie ermöglicht einen indexbasierten Zugriff auf die Elemente der Input"~, Logic- und Output-Listen sowie das Ausführen, Stoppen und Zurücksetzen einer eingegebenen Konfiguration. Dabei besteht die Möglichkeit, eine Konfiguration (\emph{PreparedConfiguration}) zunächst vollständig vorzubereiten und sie erst dann zu aktivieren, wenn sie fertig konfiguriert ist. Während der Konfigurationsphase kann noch eine andere Konfiguration (\emph{ActiveConfiguration}) aktiv sein. Beide Konfigurationen können über die \emph{ChoreographyParticipantManager}-Schnittstelle angezeigt werden, aber nur die \emph{PreparedConfiguration} kann bearbeitet werden.


\subsubsection{Konfigurierbare Kommunikation}

Abschnitt~\ref{subsec:ChoreoLeaArchitektur} stellt die Komponente \emph{OpcUaClientServerManager} zur Steuerung der Ausführung der Konfigurierbaren Kommunikation auf Basis von OPC~UA Client/ Server gemäß \cite{Stutz.2026} vor. Zur Interaktion mit dieser wird im Rahmen dieser Dissertation eine gleichnamige MTP-Schnittstellendefinition \emph{OpcUaClientServerManager}\footnote{Die \emph{OpcUaClientServerManager}-Schnittstelle kann auch als dynamisches Bedienbildelement in das LEA-Bedienbild eingebunden werden.} spezifiziert. Sie ermöglicht einen indexbasierten Zugriff auf die Listen der Verbindungen (\emph{Connections}), der Leseoperationen (\emph{Readers}), der Schreiboperationen (\emph{Writers}) sowie der schreibbaren Wertefelder (\emph{ValueFields}). Für jede Verbindung werden notwendige Variablen zur Konfiguration sowie zum Verbindungsaufbau und -abbau in Anlehnung an die \emph{UA\_Connect}- und \emph{UA\_Disconnect}-Schnittstellen der PLCopen-Spezifikation \cite{OPCFoundation.30001} bereitgestellt. Für die Lese- und Schreiboperationen werden notwendige Variablen zur Zuordnung der Operation zu einer Verbindung sowie zur Konfiguration der zu lesenden bzw. zu schreibenden Variablen auf Basis der \emph{UA\_ReadList}- bzw. \emph{UA\_WriteList}-Schnittstellen der PLCopen-Spezifikation \cite{OPCFoundation.30001} bereitgestellt. Weiterhin erfolgt eine Zuordnung jeder Leseoperation zu einem Index der Input-Liste, auf den der gelesene Wert übertragen werden soll. Jeder Schreiboperation wird ein Index der Output-Liste zugeordnet, aus dem der zu schreibende Wert ausgelesen werden soll. Für jedes Wertefeld werden Variablen zur Konfiguration des Datentyps und zur Anzeige des aktuellen Werts des Felds bereitgestellt. Für alle Konfigurationen der Verbindungen, der Lese- und der Schreiboperationen sowie für alle Wertefelder wird ähnlich wie beim \emph{ChoreographyParticipantManager} zwischen einer \emph{PreparedConfiguration} und einer \emph{ActiveConfiguration} unterschieden. 

In Abschnitt~\ref{subsec:ChoreoLeaArchitektur} wird der Datentyp \emph{UnionType} zur Übertragung von Werten mit zur Laufzeit definierbaren Datentypen gemäß \cite{Stutz.2026} vorgestellt. Zur Umsetzung dieses Datentyps wird die MTP-Schnittstelle \emph{UnionElement} eingeführt, welche die Anzeige eines Werts verschiedener Datentypen ermöglicht. Die Kommunikationsvariante des aktiven Schreibens erfordert zudem das Schreiben auf einen Wert mit definierbarem Datentyp. Dies wird durch die MTP-Schnittstellendefinition \emph{WritableUnionElement} ermöglicht. 

\begin{grayboxEnumerate}
	\textbf{Hinweis zum herstellerneutralen Informationsaustausch:}
	\begin{itemize}
		\item Das Entwurfsmuster Configurable Communication ermöglicht grundsätzlich den Austausch aller in einer LEA verfügbaren Informationen mit anderen LEAs einer choreografierten Logistics Line. Um eine herstellerneutrale Interaktion zwischen den LEAs zu ermöglichen, sollten allerdings, wo immer möglich, Variablen der standardisierten MTP-basierten LEA-Schnittstellen gemäß Kapitel~\ref{chap:Art1LEA} ausgetauscht werden.
		\item Eine Ausnahme bildet die Schnittstelle eines External Lead (siehe auch Abschnitt~\ref{subsec:LineSchnittstelle}). Hierfür werden gemäß \cite{Stutz.2026} \enquote{DoRequest-} und \enquote{DoneReply-Relationen} empfohlen, die die Zusammenführung von einem unterschiedlichen Anfang und Ende einer Choreografierelation in mehreren LEA-Diensten ermöglichen.
	\end{itemize}
\end{grayboxEnumerate}

\subsection{ChoreographySet}
\label{subsec:ChoreoSet}

% - Beschreibung der gegebenen Inhalte einer LEA (Innterne Inputs und Outputs)
% - Beschreibung der möglichen konfigurierbaren Inhalte
% - Beschreibung der Schnittstellen zur Konfiguration
% - Resultat neues AspektSet im MTP3

Zur Abbildung der in den Abschnitten~\ref{subsec:LeaChoreoModelle} und \ref{subsec:LeaChoreoSchnittstellen} eingeführten Modell- und Schnittstellendefinitionen wird ein neuer MTP-Aspekt \emph{ChoreographySet} eingeführt. Dieser Abschnitt gibt einen Überblick über diesen Aspekt, eine detaillierte Beschreibung erfolgt im Anhang~\ref{subsec:AnhangChoreographySet}. Zur Integration in einen Choreografie-Konfigurator sind die Input- und Output-Listen einer LEA mit den in Abschnitt~\ref{subsec:LeaChoreoModelle} beschriebenen Elementen modelliert. Allen Input- und Output-Elementen der Typen \emph{FixedInputElement}, \emph{FixedOutputElement}, \emph{ConfigurableInputElement} und \emph{ConfigurableOutputElement} wird eine Schnittstelle vom Typ \emph{UnionElement} zugeordnet. Jedem \emph{WritableInputElement} wird eine Schnittstelle vom Typ \emph{WritableUnionElement} zugeordnet. Die Konfiguration der Konfigurierbaren Kommunikation erfolgt mittels einer Schnittstelle vom Typ \emph{OpcUaClientServerManager}. Diese enthält für jedes \emph{ConfigurableInputElement} einen zugehörigen \emph{UaReader} und für jedes \emph{ConfigurableOutputElement} einen zugehörigen \emph{UaWriter}, die gemäß der zu ladenden Choreografiekonfiguration konfiguriert werden können. Zudem enthält diese Schnittstelle für jedes \emph{WritableInputElement} ein zugehöriges \emph{ValueField}, auf das von anderen LEAs geschrieben werden kann. Die Konfiguration der Konfigurierbaren Logik erfolgt mittels einer Schnittstelle vom Typ \emph{ChoreographyParticipantManager}. Diese ist der Modelldefinition \emph{ChoreographyParticipant} zugeordnet. Die Nutzung der neu eingeführten MTP-Schnittstellendefinitionen und -Mo\-dell\-de\-fi\-ni\-tio\-nen für die Umsetzung von Choreografierelationen wird im Anhang~\ref{subsec:AnhangChoreoRelationen} für die Prinzipien des aktiven Lesens und des aktiven Schreibens unter Nutzung von OPC~UA Client/Server beschrieben.


\section{Vertikale Integration einer Logistics Line in einen Logistics Orchestration Layer}
\label{sec:ChoreoVertikal}

Dieser Abschnitt zeigt die vertikale Integration choreografierter Logistics Lines in einen übergeordneten LOL. Hierfür wird eine Methodik zur Zusammenführung der Schnittstellen einzelner LEAs zu einer Line-Schnittstelle (Abschnitt~\ref{subsec:LineSchnittstelle}) und zur Zusammenführung der Bedienbilder einzelner LEAs zu einem Line-Bedienbild (Abschnitt~\ref{subsec:LinienBedienbild}) vorgestellt sowie die Beschreibung von Line-Schnittstellen und Line-Bedienbildern in einem Composed-MTP (Abschnitt~\ref{subsec:KonzeptComposedMTP}) behandelt.

\subsection{Line-Schnittstelle}
\label{subsec:LineSchnittstelle}

Über die Line-Schnittstelle wird die Funktionalität einer choreografierten Logistics Line einem übergeordneten LOL zur Verfügung gestellt. Die Logistics Line kann über diese Schnittstelle so parametriert, angesteuert und überwacht werden, als wäre sie ein einziger MTP-Dienst, und soll sich auch genauso verhalten. 

Die Schnittstelle eines jeden LEA-Dienstes gemäß Kapitel~\ref{chap:Art1LEA} besteht aus den in Abbildung~\ref{fig:LeaSchnittsttelle} dargestellten Bestandteilen. Genau eine \emph{ServiceControl}-Schnittstelle ist gemäß \cite{PNO.2025.Part4} verpflichtender Bestandteil eines jeden MTP-Dienstes. Sie gewährt Zugriff auf den dem Dienst zugrunde liegenden Zustandsautomaten und dessen Prozeduren. Als optionale Bestandteile enthält die Diensteschnittstelle zudem eine beliebige Anzahl an Parametern, Reportwerten und Prozesswerten. In \cite{PNO.2025.Part4} sowie in Kapitel~\ref{chap:Art1LEA} sind dafür Schnittstellen verschiedener Datentypen vorgesehen. 

\begin{figure}[htbp]
	\centerline{\includegraphics[width=0.35\textwidth]{Inhalt/Abbildungen/07_Linen-Automatisierung/LEA_Schnittstelle.pdf}}
	\caption{Bestandteile der Schnittstelle eines LEA-Dienstes}
	\label{fig:LeaSchnittsttelle}
\end{figure}

Die Line-Schnittstelle ist eine Kombination der Schnittstellen mehrerer an der Choreografie beteiligter LEA-Dienste. Dabei werden allerdings nicht alle Schnittstellenbestandteile der einzelnen LEA-Dienste auch in die Line-Schnittstelle über\-nom\-men. Ein Teil der LEA-Diensteschnittstellen wird durch die choreografieinterne Logik angesteuert und ist somit nicht von außen zugänglich. Abbildung~\ref{fig:LinienSchnittsttelle} zeigt beispielhaft, wie sich eine Line-Schnittstelle zusammensetzen kann.

\begin{figure}[htbp]
	\centerline{\includegraphics[width=0.9\textwidth]{Inhalt/Abbildungen/07_Linen-Automatisierung/Linien_Schnittstelle.pdf}}
	\caption{Bestandteile der Schnittstelle einer choreografierten Logistics Line}
	\label{fig:LinienSchnittsttelle}
\end{figure}

\subsubsection{Dienste- und Prozedurschnittstellen}

In jeder Choreografie gibt es einen führenden Dienst (in \cite{Stutz.2026} \textbf{Lead Service} genannt), welcher seinen Zustandsautomaten als Zustandsautomat der Choreografie zur Verfügung stellt. Das bedeutet, der Zustand des Lead Service spiegelt den Zustand der gesamten Choreografie wider. Zudem nimmt der Lead Service Befehle an die Choreografie (z.~B. zum Starten oder Stoppen) entgegen und ermöglicht die Einstellung einer Prozedur für die Choreografie. Als Lead Service kann einer der ohnehin an der Choreografie beteiligten LEA-Dienste genutzt werden (in \cite{Stutz.2026} \textbf{Internal Lead} genannt) oder ein separater Dienst, der ausschließlich als Lead Service in die Choreografie eingebunden wird (in \cite{Stutz.2026} \textbf{External Lead} genannt). Im Falle einer Logistics Line bietet sich die Nutzung eines External Lead an, da durch die Linienstruktur die meisten Funktionen (z.~B. Anfahren) an einer LEA initiiert werden und an einer anderen LEA enden (z.~B. Initiieren an der letzten LEA und geordnetes Anfahren bis zur ersten LEA). Ein External Lead kann dabei auf einer der LEA-Steuerungen oder auf einem externen System laufen und muss den in Abschnitt~\ref{sec:ChoreoHorizontal} beschriebenen Konzepten folgen.

Die \emph{ServiceControl}-Schnittstelle des Lead Service wird 1:1 als \emph{ServiceControl}-Bestandteil der Line-Schnittstelle übernommen (Abbildung~\ref{fig:LinienSchnittsttelle}, \enquote{ServiceControl3}). Gleichzeitig werden auch die Prozeduren des Lead Service und deren \emph{ProcedureHealthView}-Schnittstellen in die Line-Schnittstelle übernommen. Durch die Logik der Choreografie müssen die Befehle bzw. Prozedureinstellungen an der \emph{ServiceControl}-Schnittstelle so an die anderen LEA-Dienste (in \cite{Stutz.2026} \textbf{Follower Services} genannt) weitergegeben werden, dass die Choreografie ihre gewünschte Funktion erfüllt. 

Da der Lead Service durch einen Automatismus im LOL oder durch einen Operator steuerbar sein muss, wird er in der Betriebsart \emph{AutomaticExternal} oder \emph{Operator} gemäß \cite{PNO.2025.Part4} betrieben. Alle Follower Services werden ausschließlich durch LEA-interne Logik (auf Basis der Konfigurierbaren Logik) gesteuert und werden folglich in der Betriebsart \emph{AutomaticInternal} betrieben. Die Konfiguration der Betriebsarten erfolgt durch Verhaltensregeln in der Konfigurierbaren Logik.

\subsubsection{Parameter-, Reportwert-, Prozesswert-Schnittstellen}

Die Parameter, Reportwerte und Prozesswerte der Line-Schnittstelle setzen sich zusammen aus ausgewählten Parameter-, Reportwert- und Prozesswert-Schnittstellen der an der Choreografie beteiligten LEAs, welche auch für die Gesamtfunktionalität der choreografierten Logistics Line relevant sind (Abbildung~\ref{fig:LinienSchnittsttelle}, \enquote{Param1.1}, \enquote{Param1.3}, \enquote{RV2.2} und \enquote{PV3.2}). Diese Schnittstellenbestandteile werden 1:1 in die Line-Schnittstelle übernommen und können von außerhalb der Choreografie gesetzt bzw. weiterverarbeitet werden. Alle weiteren Parameter-, Reportwert-, Prozesswert-Schnittstellen werden nach Bedarf choreografieintern gesetzt und weiterverarbeitet. 

Alle Parameter, die an der Line-Schnittstelle zur Verfügung stehen, müssen durch ein automatisiertes Auftrags-Management oder einen Operator steuerbar sein. Sie werden folglich in der Betriebsart \emph{AutomaticExternal} oder \emph{Operator} gemäß \cite{PNO.2025.Part4} betrieben. Alle Parameter, die ausschließlich choreografieintern gesetzt und verarbeitet werden, werden durch LEA-interne Logik (auf Basis der Konfigurierbaren Logik) gesteuert und werden folglich in der Betriebsart \emph{Automatic\-Internal} betrieben. Die Konfiguration der Betriebsarten für Parameter wird durch Verhaltensregeln in der Konfigurierbaren Logik vorgenommen.

\subsection{Line-Bedienbild}
\label{subsec:LinienBedienbild}

Zur Anzeige und manuellen Bedienung der Logistics Line werden die Bedienbilder der LEAs zu einem Line-Bedienbild zusammengeführt. Für diese Zusammenführung sind vier verschiedene Varianten gemäß Tabelle~\ref{tab:BedienbildZusammenführung} zu unterscheiden, die auch kombiniert werden können. Zur Veranschaulichung dieser Varianten zeigt Abbildung~\ref{fig:HmiLineLeasEinzeln} die Bedienbilder von drei beispielhaften LEAs und die nachfolgenden Abbildungen~\ref{fig:HmiLineVariante1} bis \ref{fig:HmiLineVariante34} zeigen die resultierenden Line-Bedienbilder unter Nutzung der vier vorgestellten Varianten der Zusammenführung.

\begin{figure}[htbp]
	\centerline{\includegraphics[width=1.0\textwidth]{Inhalt/Abbildungen/07_Linen-Automatisierung/Line_HMI_LEAs_Einzeln.pdf}}
	\caption{Beispielhafte Bedienbilder von drei Logistics Equipment Assemblies}
	\label{fig:HmiLineLeasEinzeln}
\end{figure}

% Varianten der Bedienbildzusammenführung
\begin{footnotesize}
	\begin{longtable}{|
		>{\columncolor[HTML]{E0E0E0}}l
		>{\columncolor[HTML]{FFFFFF}}l |}
		\caption{Varianten der Zusammenführung von LEA-Bedienbildern zu Bedienbildern einer Logistics Line}
		\label{tab:BedienbildZusammenführung}\\
		\hline

		\multicolumn{1}{|L{3cm}|}{\cellcolor[HTML]{E0E0E0} \textbf{Name}}        	& 
		\multicolumn{1}{L{11.63cm}|}{\cellcolor[HTML]{E0E0E0}\textbf{Beschreibung}} 
		\\ \hline

		\multicolumn{1}{|L{3cm}|}{\cellcolor[HTML]{FFFFFF} \textbf{Variante 1:}\newline Zusammenführung von Bedienbildobjekten in einem Bedienbild}        	& 
		\multicolumn{1}{L{11.63cm}|}{\cellcolor[HTML]{FFFFFF} Diese Variante sieht eine Zusammenführung auf Ebene der Bedienbildobjekte vor. Das Bedienbild der Logistics Line enthält ausgewählte Bedienbildobjekte der einzelnen LEA-Bedienbilder. Diese werden in geeigneter Weise auf dem Line-Bedienbild angeordnet und ggf. durch neue statische Bedienbildelemente ergänzt, die Line-spezifische Gegebenheiten darstellen (z.~B. Verbindungslinien). } 
		\\ \hline

		\multicolumn{1}{|L{3cm}|}{\cellcolor[HTML]{FFFFFF} \textbf{Variante 2:}\newline Zusammenführung von Bedienbildern in einem Bedienbild}        	& 
		\multicolumn{1}{L{11.63cm}|}{\cellcolor[HTML]{FFFFFF} Diese Variante sieht eine Zusammenführung auf Ebene der Bedienbilder vor. Die Bedienbilder der einzelnen LEAs werden in das Bedienbild der Logistics Line integriert. Somit entstehen Bedienbilder, die andere Bedienbilder enthalten. }
		\\ \hline

		\multicolumn{1}{|L{3cm}|}{\cellcolor[HTML]{FFFFFF} \textbf{Variante 3:}\newline Zusammenführung von Bedienbildern in einer Hierarchie}        	& 
		\multicolumn{1}{L{11.63cm}|}{\cellcolor[HTML]{FFFFFF} Diese Variante sieht ebenfalls eine Zusammenführung auf Ebene der Bedienbilder vor. Die Bedienbilder der einzelnen LEAs werden in diesem Fall allerdings in die Bedienbildhierarchie der Logistics Line integriert. Dabei ist die Einordnung in verschiedene Hierarchieebenen möglich. }
		\\ \hline

		
		\multicolumn{1}{|L{3cm}|}{\cellcolor[HTML]{FFFFFF} \textbf{Variante 4:}\newline Zusammenführung von Bedienbildhierarchien}        	& 
		\multicolumn{1}{L{11.63cm}|}{\cellcolor[HTML]{FFFFFF} Diese Variante sieht eine Zusammenführung auf Ebene der Bedienbildhierarchien vor. Die kompletten Bedienbildhierarchien der einzelnen LEAs werden in die Bedienbildhierarchie der Logistics Line integriert. Dabei ist die Einordnung in verschiedene Hierarchieebenen möglich. }
		\\ \hline
		
	\end{longtable}
\end{footnotesize}

\begin{figure}[htbp]
	\centerline{\includegraphics[width=1.0\textwidth]{Inhalt/Abbildungen/07_Linen-Automatisierung/Line_HMI_Variante_1.pdf}}
	\caption{Zusammenführung eines Logistics-Line-Bedienbilds (Variante~1)}
	\label{fig:HmiLineVariante1}
\end{figure}

\begin{figure}[htbp]
	\centerline{\includegraphics[width=1.0\textwidth]{Inhalt/Abbildungen/07_Linen-Automatisierung/Line_HMI_Variante_2.pdf}}
	\caption{Zusammenführung eines Logistics-Line-Bedienbilds (Variante~2)}
	\label{fig:HmiLineVariante2}
\end{figure}

\begin{figure}[htbp]
	\centerline{\includegraphics[width=0.75\textwidth]{Inhalt/Abbildungen/07_Linen-Automatisierung/Line_HMI_Variante_3_4.pdf}}
	\caption{Zusammenführung eines Logistics-Line-Bedienbilds (Varianten~3 und 4)}
	\label{fig:HmiLineVariante34}
\end{figure}

\subsection{Composed-MTP}
\label{subsec:KonzeptComposedMTP}

Zur automatisierten vertikalen Integration einer choreografierten Logistics Line in einen LOL wird ein gemeinsames MTP für die Line eingeführt, welches als \textbf{Composed-MTP} bezeichnet wird. Es dient der Beschreibung der Line-Schnittstelle (Abschnitt~\ref{subsec:LineSchnittstelle}) und des Line-Bedienbilds (Abschnitt~\ref{subsec:LinienBedienbild}). Bei der Integration eines Composed-MTP werden lediglich die für die Line relevanten Informationen der LEAs in den LOL importiert und verarbeitet. Nachfolgend wird zunächst die Grundsatzentscheidung begründet, dass ein Composed-MTP Inhalte anderer MTPs referenziert, anstatt sie nachzumodellieren. Anschließend werden die Struktur, der Referenzierungsmechanismus und die Aspekte eines Composed-MTP sowie notwendige Verifizierungsworkflows vorgestellt.

\subsubsection{Nachmodellierung vs. Referenzierung}
\label{subsec:AnhangNachmodellierungReferenzierung}

Line-Schnittstelle und Line-Bedienbild, welche im Composed-MTP modelliert werden sollen, setzen sich aus Schnittstellen- und Modellelementen der an der Logistics Line beteiligten LEAs zusammen. Diese sind bereits in den LEA-MTPs modelliert und können im Zuge der Modellierung des Composed-MTP wiederverwendet werden. Hierbei bestehen zwei mögliche Varianten, die in Tabelle~\ref{tab:BewertungReferenzierungComposedMtp} beschrieben und bewertet werden.

% Bewertung Referenzierung im Composed-MTP
\begin{footnotesize}
	\begin{longtable}{|
		>{\columncolor[HTML]{E0E0E0}}l 
		>{\columncolor[HTML]{FFFFFF}}l |}
		\caption{Bewertung von Modellierungsmöglichkeiten für Composed-MTPs}
		\label{tab:BewertungReferenzierungComposedMtp}\\
		\hline
 
		\multicolumn{2}{|L{14.8cm}|}{\cellcolor[HTML]{E0E0E0}\textbf{1) Nachmodellierung}} 
		\\ \hline

		\multicolumn{2}{|L{14.8cm}|}{\cellcolor[HTML]{FFFFFF}{Die für die Line-Schnittstelle und das Line-Bedienbild relevanten Modelle und Schnittstellen der LEA-MTPs werden im Modell des Composed-MTP nachmodelliert. Die ursprünglichen LEA-MTPs werden zur Beschreibung des choreografierten Typs nicht direkt verwendet. }} 
		\\ \hline

		\multicolumn{1}{|L{1.5cm}|}{\cellcolor[HTML]{FFFFFF} Vorteile }        	& 
		\multicolumn{1}{L{13.13cm}|}{\cellcolor[HTML]{FFFFFF} Das Composed-MTP ist dem MTP einer einzelnen LEA sehr ähnlich und kann auf sehr ähnliche Weise in einen übergeordneten LOL integriert werden.} 
		\\ \hline

		\multicolumn{1}{|L{1.5cm}|}{\cellcolor[HTML]{FFFFFF} Nachteile }        	& 
		\multicolumn{1}{L{13.13cm}|}{\cellcolor[HTML]{FFFFFF} Die Modellierung des Composed-MTP doppelt sich mit der Modellierung der LEA-MTPs. Es sind Mechanismen notwendig, die sicherstellen, dass die Modellierungen in den LEA-MTPs und in den Composed-MTPs übereinstimmen. Bei Änderungen muss sowohl das LEA-MTP als auch das Composed-MTP angepasst werden.   } 
		\\ \hline

		\multicolumn{2}{|L{14.8cm}|}{\cellcolor[HTML]{E0E0E0}\textbf{2) Referenzierung}} 
		\\ \hline

		\multicolumn{2}{|L{14.8cm}|}{\cellcolor[HTML]{FFFFFF}{Die LEA-MTPs werden dem Composed-MTP zugänglich gemacht. Aus dem Modell des Composed-MTP heraus wird, wo immer möglich, auf Modelle und Schnittstellen der LEA-MTPs referenziert. Die LEA-MTPs sind somit die Basis für die Beschreibung des choreografierten Typs. }} 
		\\ \hline

		\multicolumn{1}{|L{1.5cm}|}{\cellcolor[HTML]{FFFFFF} Vorteile }        	& 
		\multicolumn{1}{L{13.13cm}|}{\cellcolor[HTML]{FFFFFF} Nur die Logistics-Line-spezifischen Inhalte werden im Composed-MTP modelliert, es erfolgt keine doppelte Modellierung im Vergleich zu den LEA-MTPs. Zudem kann auf Basis der Paketsignatur der LEA-MTPs sichergestellt werden, dass die referenzierten Modelle und Schnittstellen der LEA-MTPs nicht verändert wurden (Integrität).} 
		\\ \hline

		\multicolumn{1}{|L{1.5cm}|}{\cellcolor[HTML]{FFFFFF} Nachteile }        	& 
		\multicolumn{1}{L{13.13cm}|}{\cellcolor[HTML]{FFFFFF} Der für diese Variante notwendige Referenzierungsmechanismus ist eine Besonderheit von Composed-MTPs im Vergleich zu MTPs einzelner LEAs. Ein LOL, der LEA-MTPs integrieren kann, kann nicht automatisch auch Composed-MTPs integrieren. Es sind spezielle Integrationsmechanismen notwendig, um die Referenzen aufzulösen.} 
		\\ \hline
		
	\end{longtable}
\end{footnotesize}

\begin{graybox}
	\textbf{Designentscheidung:} Die Variante der \textbf{Referenzierung} wird als Modellierungsgrundlage für Composed-MTPs gewählt. Die heutige MTP-Spezifikation sieht nur \underline{einen} OPC~UA-Server pro MTP vor \cite{PNO.2025.Part5.1}. In einem Composed-MTP müssen Informationen aus mehreren OPC~UA-Servern mehrerer LEAs bezogen werden. Somit sind für die Integration eines Composed-MTP ohnehin andere Integrationsmechanismen als bei LEA-MTPs notwendig. Daher werden die Vorteile der Referenzierung höher bewertet. 
\end{graybox}


\subsubsection{Grundstruktur eines Composed-MTP}

Abbildung~\ref{fig:ComposedMtpGrundStruktur} zeigt die grundlegende Struktur eines Composed-MTP. 

\begin{figure}[htbp]
	\centerline{\includegraphics[width=0.7\textwidth]{Inhalt/Abbildungen/07_Linen-Automatisierung/Composed_MTP3.pdf}}
	\caption{Grundlegende Struktur eines Composed-MTP}
	\label{fig:ComposedMtpGrundStruktur}
\end{figure}

Modellierungsgrundlage bilden wie zuvor beschrieben die MTP-Dateien der an der Logistics Line beteiligten LEAs, die einen Großteil der zu modellierenden MTP-Bestandteile bereits enthalten. Um auf MTP-Bestandteile der LEA-MTPs zu referenzieren, werden diese dem Composed-MTP als Anhänge hinzugefügt (Abbildung~\ref{fig:ComposedMtpGrundStruktur}, \enquote{LEA-MTPs}). Aus diesen MTPs sind für die Beschreibung der Line-Schnittstelle und des Line-Bedienbilds notwendigen Bestandteile auszuwählen und in geeigneter Weise zu kombinieren. Hierzu dienen MTP-basierte semantische Modelle in den einzelnen Aspekten des Composed-MTP (Abbildung~\ref{fig:ComposedMtpGrundStruktur}, \enquote{Semantische Modelle}), die den semantischen Modellen eines LEA-MTPs ähneln.

\subsubsection{Referenzierung von Inhalten in den LEA-MTPs}

Die im Anhang des Composed-MTP hinterlegten LEA-MTPs werden gemäß den Mechanismen aus \cite{PNO.2025.Part1} in das \emph{AttachmentSet} des Composed-MTP eingebunden (Abbildung~\ref{fig:ComposedMtpReferenzierung}, grüne Referenz). Um auf konkrete Bestandteile der LEA-MTPs zu referenzieren, wird der in Abbildung~\ref{fig:ComposedMtpReferenzierung} dargestellte Referenzierungsmechanismus eingeführt, welcher zwei Referenzen kombiniert.

\begin{figure}[htbp]
	\centerline{\includegraphics[width=0.9\textwidth]{Inhalt/Abbildungen/07_Linen-Automatisierung/Composed_MTP_Referenzierung.pdf}}
	\caption{Referenzierungsmechanismus eines Composed-MTP}
	\label{fig:ComposedMtpReferenzierung}
\end{figure}

\textbf{1) Kontextzuordnung zu einem LEA-MTP:} Zunächst wird auf die MTP-Datei verwiesen, in welchem sich der zu referenzierende MTP-Bestandteil befindet. Hierzu wird eine neue Rollenklasse \emph{HasExternalMtpContext} eingeführt, die allen Objekten innerhalb eines Composed-MTP zugeordnet wird, die auf externe MTP-Bestandteile eines angehängten LEA-MTPs referenzieren müssen. Diese Rollenklasse nutzt den ID-Link-Mechanismus, um auf das im \emph{AttachmentSet} referenzierte LEA-MTP zu verweisen (Abbildung~\ref{fig:ComposedMtpReferenzierung}, blaue Referenz).

\textbf{2) Referenzieren des MTP-Bestandteils:} Als Grundlage für die Referenzierung eines konkreten externen MTP-Bestandteils werden die bestehenden Referenzierungsmechanismen LinkedObject und ID-Link gemäß \cite{PNO.2025.Part1} sowie der CustomSymbols-Mechanismus gemäß Abschnitt~\ref{subsec:StatHmiObjekte} verwendet. Diese Mechanismen wurden bisher zur Referenzierung von Objekten innerhalb eines MTP genutzt. In Kombination mit der zuvor beschriebenen Kontextzuordnung dienen sie nun auch zur Referenzierung von Objekten anderer MTPs (Abbildung~\ref{fig:ComposedMtpReferenzierung}, rote Referenz). Die Modellierung der LinkedObject-, ID-Link- bzw. CustomSymbols-Referenzen bleibt davon unberührt.

Dieser Referenzierungsmechanismus wird im Folgenden als \textbf{ContextReference}-Me\-cha\-nis\-mus bezeichnet. Eine detaillierte Beschreibung inkl. der Au\-to\-ma\-ti\-on\-ML-spezifischen Modellierung ist im Anhang~\ref{subsec:AnhangReferenzierungsmechanismus} zu finden.
 
\subsubsection{Aspekte eines Composed-MTP}

Ein Composed-MTP besitzt für die semantische Beschreibung der Line-Schnittstelle und des Line-Bedienbilds verschiedene Aspekte, die dem gleichen Grundaufbau folgen wie alle MTPs (Abbildung~\ref{fig:ComposedMtpAspekte}). Den Einstiegspunkt bildet eine Instanzhierarchie \emph{ModuleTypePackage}, die das \emph{Manifest} als Inhaltsverzeichnis des Composed-MTP enthält. Darunter befindet sich ein Element der Modelldefinition \emph{ComposedModuleTypePackage}, die in dieser Dissertation neu eingeführt wird. Im Gegensatz zur Modelldefinition \emph{ModuleTypePackage} gemäß \cite{PNO.2025.Part1} signalisiert diese, dass es sich um ein Composed-MTP handelt und enthält notwendige Informationen für die Typ-, Versions- und Instanzverifikation choreografierter Funktionen (siehe auch nachfolgender Abschnitt \enquote{Verifizierungsworkflows}). 

\begin{figure}[htbp]
	\centerline{\includegraphics[width=0.7\textwidth]{Inhalt/Abbildungen/07_Linen-Automatisierung/Composed_MTP_Aspekte.pdf}}
	\caption{Aspekte eines Composed-MTP}
	\label{fig:ComposedMtpAspekte}
\end{figure}

Im \emph{Manifest} werden die verschiedenen Aspekte des Composed-MTP referenziert, die in jeweils einer separaten Instanzhierarchie organisiert sind. Neben dem \emph{Manifest} ist der \emph{Attachment}-Aspekt in jedem Composed-MTP verpflichtend und enthält die angehängten LEA-MTPs sowie die Choreografiekonfiguration. Daneben können zur Beschreibung der Line-Schnittstelle und des Line-Bedienbilds der \emph{Service}-, der \emph{ProcessValue}- und der \emph{HMI}-Aspekt optional im Composed-MTP enthalten sein.

\textbf{Modellierung der Line-Schnittstelle:} Die Modellierung der Line-Schnittstelle erfolgt im \emph{Service}- und im \emph{ProcessValue}-Aspekt. Dort werden die Diensteschnittstelle des Lead Service sowie alle ausgewählten Parameter, Reportwerte und Prozesswerte semantisch modelliert. Unter Nutzung des ContextReference-Mechanismus werden die notwendigen MTP-Bestandteile der LEA-MTPs referenziert. Weitere Modell- oder Schnittstellendefinitionen sind nicht erforderlich. 

\textbf{Modellierung des Line-Bedienbilds:} Die Modellierung der Bedienbildhierarchie des Composed-MTP erfolgt im \emph{HMI}-Aspekt. Hierbei ist zwischen den in Abschnitt~\ref{subsec:LinienBedienbild} beschriebenen Varianten der Zusammenführung von Bedienbildern zu unterscheiden. Die Variante 1 (\emph{Zusammenführung von Bedienbildobjekten}) erfolgt unter Nutzung des ContextReference-Mechanismus und erfordert keine weiteren Modell- oder Schnittstellendefinitionen. Für die Variante 2 (\emph{Zusammenführung von Bedienbildern in einem Bedienbild}) wird eine neue Modelldefinition \emph{PictureFrame} eingeführt, welche den ContextReference-Mechanismus zur Referenzierung ganzer Bedienbilder nutzt und diese in einer bestimmten Größe und Ausrichtung auf einem anderen Bedienbild platziert. Für die Varianten 3 und 4 (\emph{Zusammenführung von Bedienbildern in einer Hierarchie} und \emph{Zusammenführung von Bedienbildhierarchien}) wird eine neue Modelldefinition \emph{ReferencedPicture} eingeführt, welche mittels des ContextReference-Mechanismus Bedienbilder bzw. Bedienbildhierarchien referenziert und diese in die Bedienbildhierarchie des Composed-MTP einfügt. 

Die AutomationML-spezifische Modellierung der zuvor genannten Aspekte eines Composed-MTP ist im Anhang~\ref{subsec:AnhangComposedMtpAspekte} beschrieben. Weitere Aspekte, z.~B. ein \emph{Text}-Aspekt zur Anzeige Line-spezifischer Texte, ein \emph{Alarm}-Aspekt für Line-spezifische Alarme und Events oder ein \emph{Diagnostics}-Aspekt für Diagnoseinformationen, können zukünftig ergänzt werden. Der \emph{ServerAssembly}-Aspekt und der \emph{DataAssembly}-Aspekt sind im Kontext choreografierter Funktionen nicht im Composed-MTP notwendig, da die betreffenden Informationen bereits in den LEA-MTPs enthalten sind und mittels des ContextReference-Mechanismus referenziert werden.

\subsubsection{Verifizierungsworkflows}

In \cite{PNO.2025.Part1} werden Typ-, Versions- und Instanzverifikationsworkflows als Workflows des \emph{Manifest}-Aspekts beschrieben. Im Rahmen dieser Dissertation werden diese Workflows um die Verifizierung choreografierter Funktionen und Composed-MTPs erweitert. 

Für die Typ- und Versionsverifikation werden \textbf{Typ- und Versionsinformationen} der choreografierten Funktion im Composed-MTP hinterlegt und können anschließend im Rahmen der Verifikation mit der tatsächlich auf die LEAs einer Logistics Line geladenen Konfiguration verglichen werden. So wird sichergestellt, dass \textbf{die richtige Choreografiekonfiguration geladen ist}.

Im Zuge der \textbf{Instanzverifikation} wird für jede notwendige Teilnehmerrolle, welche in einer choreografierten Logistics Line zu besetzen ist, verglichen, ob die in der Planung vorgesehene LEA-Instanz mit der tatsächlich eingebauten LEA-Instanz übereinstimmt. Als Grundlage für diese Verifikation werden die Instanzkennung der LEAs und eine ID für jede Teilnehmerrolle verwendet, die im Composed-MTP festgelegt ist. So wird sichergestellt, dass \textbf{die richtigen LEA-Instanzen in der Logistics Line verwendet werden}.

\begin{graybox}
	\textbf{Hinweis:} Die MTP-basierten Verifikationsworkflows prüfen \underline{\textbf{nicht}}, ob eine Choreografie ausführbar ist, oder mit welchen Parametern sie funktioniert. Sie stellen lediglich sicher, dass die richtige Choreografiekonfiguration geladen ist und die richtigen LEA-Instanzen in der Logistics Line verwendet werden.
\end{graybox}

 
\section{MTP-Erweiterungen}
\label{sec:MTPErweiterungenLL}

In den vorangegangenen Abschnitten wurde eine Reihe von Erweiterungen der MTP-Spezifikation identifiziert, die im Kontext der Automatisierung von Logistics Lines notwendig sind. Tabelle~\ref{tab:MtpErweiterungenLL} gibt einen Überblick über die hierfür einzuführenden Modell- und Schnittstellendefinitionen sowie Workflows. Eine detailliertere Beschreibung erfolgt in den in Tabelle~\ref{tab:MtpErweiterungenLL} angegebenen Anhängen.

% MTP-Erweiterungen LL
\begin{footnotesize}
	\begin{longtable}{|
		>{\columncolor[HTML]{E0E0E0}}l 
		>{\columncolor[HTML]{FFFFFF}}l 
		>{\columncolor[HTML]{FFFFFF}}l 
		>{\columncolor[HTML]{FFFFFF}}l |}
		\caption{Erweiterungen der MTP-Spezifikation für die Automatisierung von Logistics Lines}
		\label{tab:MtpErweiterungenLL}\\
		\hline

		% Schnittstellendefinitionen
		\multicolumn{4}{|L{14.8cm}|}{\cellcolor[HTML]{E0E0E0}\textbf{Schnittstellendefinitionen}} 
		\\ \hline

		\multicolumn{1}{|L{3.6cm}|}{\cellcolor[HTML]{E0E0E0}\textbf{Profil}}        	& 
		\multicolumn{1}{L{2.8cm}|}{\cellcolor[HTML]{E0E0E0}\textbf{Definition}} 			& 
		\multicolumn{1}{L{5.43cm}|}{\cellcolor[HTML]{E0E0E0}\textbf{Beschreibung}} 			& 
		\multicolumn{1}{L{2.45cm}|}{\cellcolor[HTML]{E0E0E0}\textbf{Detail-beschreibung}} 
		\\ \hline

		\multicolumn{1}{|L{3.6cm}|}{\cellcolor[HTML]{FFFFFF} ModuleTypePackage: ChoreographySet.Base V2.0.0 }                    & 
		\multicolumn{1}{L{2.8cm}|}{\cellcolor[HTML]{FFFFFF}  SUC Union\-Element}              &      
		\multicolumn{1}{L{5.43cm}|}{\cellcolor[HTML]{FFFFFF} Schnittstelle zum Auslesen eines Werts mit zur Laufzeit auswählbarem Datentyp}              &     
		\multicolumn{1}{L{2.45cm}|}{\cellcolor[HTML]{FFFFFF} Anhang~\ref{subsec:AnhangChoreographySet}}             
		\\ \hline

		\multicolumn{1}{|L{3.6cm}|}{\cellcolor[HTML]{FFFFFF} ModuleTypePackage: ChoreographySet.Base V2.0.0 }                    & 
		\multicolumn{1}{L{2.8cm}|}{\cellcolor[HTML]{FFFFFF} SUC Writable\-Union\-Element}              &      
		\multicolumn{1}{L{5.43cm}|}{\cellcolor[HTML]{FFFFFF} Schnittstelle zum Schreiben eines Werts mit zur Laufzeit auswählbarem Datentyp}              &     
		\multicolumn{1}{L{2.45cm}|}{\cellcolor[HTML]{FFFFFF} Anhang~\ref{subsec:AnhangChoreographySet}}             
		\\ \hline

		\multicolumn{1}{|L{3.6cm}|}{\cellcolor[HTML]{FFFFFF} ModuleTypePackage: ChoreographySet.Base V2.0.0 }                    & 
		\multicolumn{1}{L{2.8cm}|}{\cellcolor[HTML]{FFFFFF} SUC Choreography\-Element}              &      
		\multicolumn{1}{L{5.43cm}|}{\cellcolor[HTML]{FFFFFF} Basis-Schnittstelle für alle zur Choreografiekonfiguration notwendigen Schnittstellen}              &     
		\multicolumn{1}{L{2.45cm}|}{\cellcolor[HTML]{FFFFFF} Anhang~\ref{subsec:AnhangChoreographySet}}             
		\\ \hline

		\multicolumn{1}{|L{3.6cm}|}{\cellcolor[HTML]{FFFFFF} ModuleTypePackage: ChoreographySet.Base V2.0.0 }                    & 
		\multicolumn{1}{L{2.8cm}|}{\cellcolor[HTML]{FFFFFF} SUC Choreography\-Participant\-Manager}              &      
		\multicolumn{1}{L{5.43cm}|}{\cellcolor[HTML]{FFFFFF} Schnittstelle zur Konfiguration und Steuerung der Ausführung der Konfigurierbaren Logik einer Choreografie}              &     
		\multicolumn{1}{L{2.45cm}|}{\cellcolor[HTML]{FFFFFF} Anhang~\ref{subsec:AnhangChoreographySet}}             
		\\ \hline

		\multicolumn{1}{|L{3.6cm}|}{\cellcolor[HTML]{FFFFFF} ModuleTypePackage: ChoreographySet.Base V2.0.0 }                    & 
		\multicolumn{1}{L{2.8cm}|}{\cellcolor[HTML]{FFFFFF} SUC Communication\-Manager}              &      
		\multicolumn{1}{L{5.43cm}|}{\cellcolor[HTML]{FFFFFF} Basis-Schnittstelle für alle Konfigurationsschnittstellen der Konfigurierbaren Kommunikation einer Choreografie}              &     
		\multicolumn{1}{L{2.45cm}|}{\cellcolor[HTML]{FFFFFF} Anhang~\ref{subsec:AnhangChoreographySet}}             
		\\ \hline

		\multicolumn{1}{|L{3.6cm}|}{\cellcolor[HTML]{FFFFFF} ModuleTypePackage: ChoreographySet.Base V2.0.0 }                    & 
		\multicolumn{1}{L{2.8cm}|}{\cellcolor[HTML]{FFFFFF} SUC OpcUa\-Client\-Server\-Manager}              &      
		\multicolumn{1}{L{5.43cm}|}{\cellcolor[HTML]{FFFFFF} Schnittstelle zur Konfiguration und Steuerung der Ausführung der Konfigurierbaren Kommunikation einer Choreografie auf Basis von OPC~UA Client/Server}              &     
		\multicolumn{1}{L{2.45cm}|}{\cellcolor[HTML]{FFFFFF} Anhang~\ref{subsec:AnhangChoreographySet}}             
		\\ \hline
		

		% Modelldefinitionen
		\multicolumn{4}{|L{14.8cm}|}{\cellcolor[HTML]{E0E0E0}\textbf{Modelldefinitionen}} 
		\\ \hline

		\multicolumn{1}{|L{3.6cm}|}{\cellcolor[HTML]{E0E0E0}\textbf{Profil}}        	& 
		\multicolumn{1}{L{2.8cm}|}{\cellcolor[HTML]{E0E0E0}\textbf{Definition}} 			& 
		\multicolumn{1}{L{5.43cm}|}{\cellcolor[HTML]{E0E0E0}\textbf{Beschreibung}} 			& 
		\multicolumn{1}{L{2.45cm}|}{\cellcolor[HTML]{E0E0E0}\textbf{Detail-beschreibung}} 
		\\ \hline

		\multicolumn{1}{|L{3.6cm}|}{\cellcolor[HTML]{FFFFFF} ModuleTypePackage: ChoreographySet.Base V2.0.0 }                    & 
		\multicolumn{1}{L{2.8cm}|}{\cellcolor[HTML]{FFFFFF} IH Choreography}              &      
		\multicolumn{1}{L{5.43cm}|}{\cellcolor[HTML]{FFFFFF} Instanzhierarchie zur Verwaltung aller choreografiebezogenen Modelle eines MTP}              &     
		\multicolumn{1}{L{2.45cm}|}{\cellcolor[HTML]{FFFFFF} Anhang~\ref{subsec:AnhangChoreographySet}}             
		\\ \hline

		\multicolumn{1}{|L{3.6cm}|}{\cellcolor[HTML]{FFFFFF} ModuleTypePackage: ChoreographySet.Base V2.0.0 }                    & 
		\multicolumn{1}{L{2.8cm}|}{\cellcolor[HTML]{FFFFFF} SUC Choreography\-Set}              &      
		\multicolumn{1}{L{5.43cm}|}{\cellcolor[HTML]{FFFFFF} AspektSet, das alle choreografiebezogenen Modelldefinitionen organisiert}              &     
		\multicolumn{1}{L{2.45cm}|}{\cellcolor[HTML]{FFFFFF} Anhang~\ref{subsec:AnhangChoreographySet}}             
		\\ \hline

		\multicolumn{1}{|L{3.6cm}|}{\cellcolor[HTML]{FFFFFF} ModuleTypePackage: ChoreographySet.Base V2.0.0 }                    & 
		\multicolumn{1}{L{2.8cm}|}{\cellcolor[HTML]{FFFFFF} SUC Choreography\-Participant}              &      
		\multicolumn{1}{L{5.43cm}|}{\cellcolor[HTML]{FFFFFF} Modell zur Beschreibung einer LEA als Choreografieteilnehmer}              &     
		\multicolumn{1}{L{2.45cm}|}{\cellcolor[HTML]{FFFFFF} Anhang~\ref{subsec:AnhangChoreographySet}}             
		\\ \hline

		\multicolumn{1}{|L{3.6cm}|}{\cellcolor[HTML]{FFFFFF} ModuleTypePackage: ChoreographySet.Base V2.0.0 }                    & 
		\multicolumn{1}{L{2.8cm}|}{\cellcolor[HTML]{FFFFFF} SUC Input\-List}              &      
		\multicolumn{1}{L{5.43cm}|}{\cellcolor[HTML]{FFFFFF} Modell für die Input-Liste einer LEA}              &     
		\multicolumn{1}{L{2.45cm}|}{\cellcolor[HTML]{FFFFFF} Anhang~\ref{subsec:AnhangChoreographySet}}             
		\\ \hline

		\multicolumn{1}{|L{3.6cm}|}{\cellcolor[HTML]{FFFFFF} ModuleTypePackage: ChoreographySet.Base V2.0.0 }                    & 
		\multicolumn{1}{L{2.8cm}|}{\cellcolor[HTML]{FFFFFF} SUC Output\-List}              &      
		\multicolumn{1}{L{5.43cm}|}{\cellcolor[HTML]{FFFFFF} Modell für die Output-Liste einer LEA}              &     
		\multicolumn{1}{L{2.45cm}|}{\cellcolor[HTML]{FFFFFF} Anhang~\ref{subsec:AnhangChoreographySet}}             
		\\ \hline

		\multicolumn{1}{|L{3.6cm}|}{\cellcolor[HTML]{FFFFFF} ModuleTypePackage: ChoreographySet.Base V2.0.0 }                    & 
		\multicolumn{1}{L{2.8cm}|}{\cellcolor[HTML]{FFFFFF} SUC Input\-Element}              &      
		\multicolumn{1}{L{5.43cm}|}{\cellcolor[HTML]{FFFFFF} Modell eines Elements der Input-Liste einer LEA}              &     
		\multicolumn{1}{L{2.45cm}|}{\cellcolor[HTML]{FFFFFF} Anhang~\ref{subsec:AnhangChoreographySet}}             
		\\ \hline

		\multicolumn{1}{|L{3.6cm}|}{\cellcolor[HTML]{FFFFFF} ModuleTypePackage: ChoreographySet.Base V2.0.0 }                    & 
		\multicolumn{1}{L{2.8cm}|}{\cellcolor[HTML]{FFFFFF} SUC Output\-Element}              &      
		\multicolumn{1}{L{5.43cm}|}{\cellcolor[HTML]{FFFFFF} Modell eines Elements der Output-Liste einer LEA}              &     
		\multicolumn{1}{L{2.45cm}|}{\cellcolor[HTML]{FFFFFF} Anhang~\ref{subsec:AnhangChoreographySet}}             
		\\ \hline

		\multicolumn{1}{|L{3.6cm}|}{\cellcolor[HTML]{FFFFFF} ModuleTypePackage: ChoreographySet.Base V2.0.0 }                    & 
		\multicolumn{1}{L{2.8cm}|}{\cellcolor[HTML]{FFFFFF} SUC Fixed\-Input\-Element}              &      
		\multicolumn{1}{L{5.43cm}|}{\cellcolor[HTML]{FFFFFF} Modell eines durch die LEA-Programmierung fest vorgegebenen Elements der Input-Liste}              &     
		\multicolumn{1}{L{2.45cm}|}{\cellcolor[HTML]{FFFFFF} Anhang~\ref{subsec:AnhangChoreographySet}}             
		\\ \hline

		\multicolumn{1}{|L{3.6cm}|}{\cellcolor[HTML]{FFFFFF} ModuleTypePackage: ChoreographySet.Base V2.0.0 }                    & 
		\multicolumn{1}{L{2.8cm}|}{\cellcolor[HTML]{FFFFFF} SUC Fixed\-Output\-Element}              &      
		\multicolumn{1}{L{5.43cm}|}{\cellcolor[HTML]{FFFFFF} Modell eines durch die LEA-Programmierung fest vorgegebenen Elements der Output-Liste}              &     
		\multicolumn{1}{L{2.45cm}|}{\cellcolor[HTML]{FFFFFF} Anhang~\ref{subsec:AnhangChoreographySet}}             
		\\ \hline
		
		\multicolumn{1}{|L{3.6cm}|}{\cellcolor[HTML]{FFFFFF} ModuleTypePackage: ChoreographySet.Base V2.0.0 }                    & 
		\multicolumn{1}{L{2.8cm}|}{\cellcolor[HTML]{FFFFFF} SUC Configurable\-Input\-Element}              &      
		\multicolumn{1}{L{5.43cm}|}{\cellcolor[HTML]{FFFFFF} Modell eines Elements der Input-Liste zum Einlesen eines externen Werts aus einer anderen LEA}              &     
		\multicolumn{1}{L{2.45cm}|}{\cellcolor[HTML]{FFFFFF} Anhang~\ref{subsec:AnhangChoreographySet}}             
		\\ \hline

		\multicolumn{1}{|L{3.6cm}|}{\cellcolor[HTML]{FFFFFF} ModuleTypePackage: ChoreographySet.Base V2.0.0 }                    & 
		\multicolumn{1}{L{2.8cm}|}{\cellcolor[HTML]{FFFFFF} SUC Configurable\-Output\-Element}              &      
		\multicolumn{1}{L{5.43cm}|}{\cellcolor[HTML]{FFFFFF} Modell eines Elements der Output-Liste zum Schreiben eines Werts an eine andere LEA}              &     
		\multicolumn{1}{L{2.45cm}|}{\cellcolor[HTML]{FFFFFF} Anhang~\ref{subsec:AnhangChoreographySet}}             
		\\ \hline

		\multicolumn{1}{|L{3.6cm}|}{\cellcolor[HTML]{FFFFFF} ModuleTypePackage: ChoreographySet.Base V2.0.0 }                    & 
		\multicolumn{1}{L{2.8cm}|}{\cellcolor[HTML]{FFFFFF} SUC Writable\-Input\-Element}              &      
		\multicolumn{1}{L{5.43cm}|}{\cellcolor[HTML]{FFFFFF} Modell eines passiven Elements der Input-Liste, auf das durch eine andere LEA geschrieben werden kann}              &     
		\multicolumn{1}{L{2.45cm}|}{\cellcolor[HTML]{FFFFFF} Anhang~\ref{subsec:AnhangChoreographySet}}             
		\\ \hline

		\multicolumn{1}{|L{3.6cm}|}{\cellcolor[HTML]{FFFFFF} ModuleTypePackage: Manifest.Composed V2.0.0 }                    & 
		\multicolumn{1}{L{2.8cm}|}{\cellcolor[HTML]{FFFFFF} SUC Composed\-Module\-Type\-Package}              &      
		\multicolumn{1}{L{5.43cm}|}{\cellcolor[HTML]{FFFFFF} Basis-Modell für ein Composed-MTP }              &     
		\multicolumn{1}{L{2.45cm}|}{\cellcolor[HTML]{FFFFFF} Anhang~\ref{subsec:AnhangManifest}}             
		\\ \hline

		\multicolumn{1}{|L{3.6cm}|}{\cellcolor[HTML]{FFFFFF} ModuleTypePackage: Manifest.Composed V2.0.0 }                    & 
		\multicolumn{1}{L{2.8cm}|}{\cellcolor[HTML]{FFFFFF} AT Composed\-Type\-Revision\-Type}              &      
		\multicolumn{1}{L{5.43cm}|}{\cellcolor[HTML]{FFFFFF} Attributtyp für die Versionsinformationen einer choreografierten Funktion }              &     
		\multicolumn{1}{L{2.45cm}|}{\cellcolor[HTML]{FFFFFF} Anhang~\ref{subsec:AnhangManifest}}             
		\\ \hline

		\multicolumn{1}{|L{3.6cm}|}{\cellcolor[HTML]{FFFFFF} ModuleTypePackage: Manifest.Composed V2.0.0 }                    & 
		\multicolumn{1}{L{2.8cm}|}{\cellcolor[HTML]{FFFFFF} RC Has\-External\-Mtp\-Context}              &      
		\multicolumn{1}{L{5.43cm}|}{\cellcolor[HTML]{FFFFFF} RoleClass zur Kennzeichnung, dass ein referenziertes Objekt aus einer anderen angehängten MTP-Datei stammt }              &     
		\multicolumn{1}{L{2.45cm}|}{\cellcolor[HTML]{FFFFFF} Anhang~\ref{subsec:AnhangManifest}}             
		\\ \hline

		\multicolumn{1}{|L{3.6cm}|}{\cellcolor[HTML]{FFFFFF} ModuleTypePackage: HMISet.Composed V2.0.0 }                    & 
		\multicolumn{1}{L{2.8cm}|}{\cellcolor[HTML]{FFFFFF} SUC Picture\-Frame}              &      
		\multicolumn{1}{L{5.43cm}|}{\cellcolor[HTML]{FFFFFF} Modell zur Referenzierung von in anderen MTP-Dateien befindlichen Bedienbildern zur Nutzung in einem Bedienbild (Bild in Bild) }              &     
		\multicolumn{1}{L{2.45cm}|}{\cellcolor[HTML]{FFFFFF} Anhang~\ref{subsec:AnhangHMISet}}             
		\\ \hline

		\multicolumn{1}{|L{3.6cm}|}{\cellcolor[HTML]{FFFFFF} ModuleTypePackage: HMISet.Composed V2.0.0 }                    & 
		\multicolumn{1}{L{2.8cm}|}{\cellcolor[HTML]{FFFFFF} SUC Picture\newline (Erweiterung)}              &      
		\multicolumn{1}{L{5.43cm}|}{\cellcolor[HTML]{FFFFFF} Erweiterung des bestehenden Picture-Modells um die Möglichkeit, PictureFrames einzufügen}              &     
		\multicolumn{1}{L{2.45cm}|}{\cellcolor[HTML]{FFFFFF} Anhang~\ref{subsec:AnhangHMISet}}             
		\\ \hline

		\multicolumn{1}{|L{3.6cm}|}{\cellcolor[HTML]{FFFFFF} ModuleTypePackage: HMISet.Composed V2.0.0 }                    & 
		\multicolumn{1}{L{2.8cm}|}{\cellcolor[HTML]{FFFFFF} SUC SemanticGroup\newline (Erweiterung)}              &      
		\multicolumn{1}{L{5.43cm}|}{\cellcolor[HTML]{FFFFFF} Erweiterung des bestehenden SemanticGroup-Modells um die Möglichkeit, PictureFrames einzufügen}              &
		\multicolumn{1}{L{2.45cm}|}{\cellcolor[HTML]{FFFFFF} Anhang~\ref{subsec:AnhangHMISet}}             
		\\ \hline
		
		\multicolumn{1}{|L{3.6cm}|}{\cellcolor[HTML]{FFFFFF} ModuleTypePackage: HMISet.Composed V2.0.0 }                    & 
		\multicolumn{1}{L{2.8cm}|}{\cellcolor[HTML]{FFFFFF} SUC Referenced\-Picture}              &      
		\multicolumn{1}{L{5.43cm}|}{\cellcolor[HTML]{FFFFFF} Modell zur Referenzierung von in anderen MTP-Dateien befindlichen Bedienbildern zur Nutzung in einer Bedienbildhierarchie }              &     
		\multicolumn{1}{L{2.45cm}|}{\cellcolor[HTML]{FFFFFF} Anhang~\ref{subsec:AnhangHMISet}}             
		\\ \hline

		\multicolumn{1}{|L{3.6cm}|}{\cellcolor[HTML]{FFFFFF} ModuleTypePackage: HMISet.Composed V2.0.0 }                    & 
		\multicolumn{1}{L{2.8cm}|}{\cellcolor[HTML]{FFFFFF} SUC HMI\-Set \newline (Erweiterung)}              &      
		\multicolumn{1}{L{5.43cm}|}{\cellcolor[HTML]{FFFFFF} Erweiterung des bestehenden HMISet-Modells um die Möglichkeit, ReferencedPictures einzufügen}              &
		\multicolumn{1}{L{2.45cm}|}{\cellcolor[HTML]{FFFFFF} Anhang~\ref{subsec:AnhangHMISet}}             
		\\ \hline


		% Workflows
		\multicolumn{4}{|L{14.8cm}|}{\cellcolor[HTML]{E0E0E0}\textbf{Workflows}} 
		\\ \hline

		\multicolumn{1}{|L{3.6cm}|}{\cellcolor[HTML]{E0E0E0}\textbf{Profil}}        	& 
		\multicolumn{1}{L{2.8cm}|}{\cellcolor[HTML]{E0E0E0}\textbf{Workflow}} 			& 
		\multicolumn{1}{L{5.43cm}|}{\cellcolor[HTML]{E0E0E0}\textbf{Beschreibung}} 			& 
		\multicolumn{1}{L{2.45cm}|}{\cellcolor[HTML]{E0E0E0}\textbf{Detail-beschreibung}} 
		\\ \hline

		\multicolumn{1}{|L{3.6cm}|}{\cellcolor[HTML]{FFFFFF} ModuleTypePackage: Manifest.Composed V2.0.0 }                    & 
		\multicolumn{1}{L{2.8cm}|}{\cellcolor[HTML]{FFFFFF} Typverifikation}              &      
		\multicolumn{1}{L{5.43cm}|}{\cellcolor[HTML]{FFFFFF} Workflow zur Verifikation des Typs einer auf die LEAs geladenen choreografierten Funktion gegen den in einem Composed-MTP beschriebenen Typ }              &     
		\multicolumn{1}{L{2.45cm}|}{\cellcolor[HTML]{FFFFFF} Anhang~\ref{subsec:AnhangManifest}}             
		\\ \hline

		\multicolumn{1}{|L{3.6cm}|}{\cellcolor[HTML]{FFFFFF} ModuleTypePackage: Manifest.Composed V2.0.0 }                    & 
		\multicolumn{1}{L{2.8cm}|}{\cellcolor[HTML]{FFFFFF} Versions-verifikation}              &      
		\multicolumn{1}{L{5.43cm}|}{\cellcolor[HTML]{FFFFFF} Workflow zur Verifikation der Version einer auf die LEAs geladenen choreografierten Funktion gegen die in einem Composed-MTP beschriebene Version }              &  
		\multicolumn{1}{L{2.45cm}|}{\cellcolor[HTML]{FFFFFF} Anhang~\ref{subsec:AnhangManifest}}             
		\\ \hline

		\multicolumn{1}{|L{3.6cm}|}{\cellcolor[HTML]{FFFFFF} ModuleTypePackage: Manifest.Composed V2.0.0 }                    & 
		\multicolumn{1}{L{2.8cm}|}{\cellcolor[HTML]{FFFFFF} Instanz-verifikation}              &      
		\multicolumn{1}{L{5.43cm}|}{\cellcolor[HTML]{FFFFFF} Workflow zur Verifikation der für eine choreografierte Funktion eingesetzten LEA-Instanzen gegen die geplanten LEA-Instanzen  }              &     
		\multicolumn{1}{L{2.45cm}|}{\cellcolor[HTML]{FFFFFF} Anhang~\ref{subsec:AnhangManifest}}             
		\\ \hline

	\end{longtable}
\end{footnotesize}

Für die choreografiebezogenen Modelldefinitionen wird zudem eine neue Bibliothek (\emph{SUCL MTPChoreographySUCLib}) eingeführt (Anhang~\ref{subsec:AnhangChoreographySet}). Für die \emph{RC HasExternalMtpContext} wird eine neue Bibliothek (\emph{RCL MTPRCLib}) eingeführt (Anhang~\ref{subsec:AnhangManifest}).